
\documentclass{article}
\usepackage{colortbl}
\usepackage{makecell}
\usepackage{multirow}
\usepackage{supertabular}

\begin{document}

\newcounter{utterance}

\centering \large Interaction Transcript for game `eqbench', experiment `eq\_bench', episode 18 with reflexion{-}Qwen3.5{-}27B.
\vspace{24pt}

{ \footnotesize  \setcounter{utterance}{1}
\setlength{\tabcolsep}{0pt}
\begin{supertabular}{c@{$\;$}|p{.15\linewidth}@{}p{.15\linewidth}p{.15\linewidth}p{.15\linewidth}p{.15\linewidth}p{.15\linewidth}}
    \# & \multicolumn{2}{c}{Player} && \multicolumn{2}{c}{Game Master} \\
    \hline

    \theutterance \stepcounter{utterance}  
    & & & \multicolumn{4}{p{0.6\linewidth}}{
        \cellcolor[rgb]{0.9,0.9,0.9}{
            \makecell[{{p{\linewidth}}}]{
                \texttt{\tiny{[P1$\langle$GM]}}
                \texttt{You have attempted this same task before and did not succeed. These are the plans you wrote after those attempts.} \\
\texttt{Trial \#1: New plan: My previous attempt failed because I assigned a non{-}zero score to "Angry" (4) and "Relieved" (2), which likely contradicted the specific emotional nuance of the character's final line ("I'm done here"). Dr. Sinclair's tone is one of exhausted dismissal and intellectual exasperation, not active rage or relief. The primary emotion is clearly "Frustrated" due to the futility of arguing with bad faith. "Sympathetic" should be 0 as he explicitly calls the other a "sensationalist." To correct this, I will assign a high score (e.g., 9 or 10) to "Frustrated" and strictly assign 0 to "Angry," "Relieved," and "Sympathetic," ensuring the scores reflect the singular dominant emotion of the scene without diluting it with secondary, unsupported feelings.} \\
\\ 
\texttt{Trial \#2: New plan: My previous attempt failed because I assigned a score of 0 to "Angry," overlooking the aggressive and confrontational nature of Dr. Sinclair's final lines ("You're living in a nightmare," "I'm done here"). While "Frustrated" is the dominant emotion, the character's tone has escalated from exasperation to active hostility and anger at the end of the dialogue. To correct this, I will assign a high score (e.g., 9) to "Frustrated" as the primary emotion, but also assign a significant non{-}zero score (e.g., 7) to "Angry" to reflect the sharp, cutting nature of his final rebuttal. I will maintain "Relieved" and "Sympathetic" at 0, as the character shows no relief or pity, only dismissal and irritation.} \\
\\ 
\texttt{Trial \#3: New plan: My previous attempt failed because I assigned a non{-}zero score to "Angry" (7), which likely misinterpreted Dr. Sinclair's final tone. While his words are sharp ("nightmare," "I'm done here"), the context of a long, futile debate with a conspiracy theorist suggests his primary state is deep exasperation and exhaustion rather than active rage. The phrase "I'm done here" signals a desire to escape the interaction, not an outburst of anger. To correct this, I will assign the maximum score (10) to "Frustrated" to capture the intensity of his exasperation, while strictly assigning 0 to "Angry," "Relieved," and "Sympathetic." This ensures the output reflects the singular, dominant emotion of intellectual exhaustion and dismissal without diluting it with the secondary emotion of anger.} \\
\\ 
\\ 
\\ 
\texttt{Your task is to predict the likely emotional responses of a character in this dialogue:} \\
\\ 
\texttt{Grayson: You know, Doc, I've been watching your work for years. All that time and money spent on Mars, and what have you got to show for it? } \\
\texttt{Dr. Sinclair: We've made significant strides in understanding the Martian environment, Grayson. It's not about instant results. } \\
\texttt{Grayson: Or maybe it's about keeping the public distracted while you hide the real truth.} \\
\texttt{Dr. Sinclair: And what truth would that be, Mr. Hawk?} \\
\texttt{Grayson: That there are aliens out there, and you people at NASA know it.} \\
\texttt{Dr. Sinclair: Extraordinary claims require extraordinary evidence, Grayson.} \\
\texttt{Grayson: Oh, I have evidence. I've got testimonials, photos, documents...} \\
\texttt{Dr. Sinclair: Anecdotes and blurry photos aren't evidence, Grayson. They're stories. } \\
\texttt{Grayson: Stories that you're too afraid to believe.} \\
\texttt{Dr. Sinclair: No, stories that don't hold up under scientific scrutiny. } \\
\texttt{Grayson: You're just a puppet, Sinclair. They've got you so convinced of your 'facts' that you can't see the truth.} \\
\texttt{Dr. Sinclair: And you're a sensationalist, Hawk. You've built a career on paranoia and misinformation.} \\
\texttt{Grayson: At least I'm not lying to the public.} \\
\texttt{Dr. Sinclair: I'm not lying, Grayson. I'm presenting facts, not theories.} \\
\texttt{Grayson: Your 'facts' are just what they want you to believe. } \\
\texttt{Dr. Sinclair: And your theories are what your followers want to hear. It's a comforting narrative, isn't it? The idea that there's some grand conspiracy, that you're the hero fighting against it.} \\
\texttt{Grayson: And you're just a cog in the machine, Sinclair. You're part of the problem.} \\
\texttt{Dr. Sinclair: I'm part of the solution, Grayson. I'm trying to educate people, not feed them fantasies.} \\
\texttt{Grayson: You're the fantasy, Sinclair. You're living in a dream world.} \\
\texttt{Dr. Sinclair: And you're living in a nightmare of your own making. I'm done here.} \\
\texttt{{[}End dialogue{]}} \\
\\ 
\texttt{At the end of this dialogue, Dr. Sinclair would feel...} \\
\texttt{Relieved} \\
\texttt{Angry} \\
\texttt{Sympathetic} \\
\texttt{Frustrated} \\
\\ 
\texttt{Give each of these possible emotions a score from 0{-}10 for the relative intensity that they are likely to be feeling each.} \\
\\ 
\texttt{You must output in the following format, including headings (of course, you should give your own scores), with no additional commentary:} \\
\\ 
\texttt{Relieved: <score>} \\
\texttt{Angry: <score>} \\
\texttt{Sympathetic: <score>} \\
\texttt{Frustrated: <score>} \\
\\ 
\\ 
\texttt{{[}End of answer{]}} \\
\\ 
\texttt{Remember: zero is a valid score, meaning they are likely not feeling that emotion. You must score at least one emotion > 0.} \\
\\ 
\texttt{Your answer:} \\
            }
        }
    }
     \\ \\

    \theutterance \stepcounter{utterance}  
    & \multicolumn{4}{p{0.6\linewidth}}{
        \cellcolor[rgb]{0.561,0.737,0.561}{
            \makecell[{{p{\linewidth}}}]{
                \texttt{\tiny{[P1$\rangle$GM]}}
                \texttt{Relieved: 0} \\
\texttt{Angry: 0} \\
\texttt{Sympathetic: 0} \\
\texttt{Frustrated: 10} \\
            }
        }
    }
    & & \\ \\

    \theutterance \stepcounter{utterance}  
    & & & \multicolumn{2}{p{0.3\linewidth}}{
        \cellcolor[rgb]{0.9,0.9,0.9}{
            \makecell[{{p{\linewidth}}}]{
                \texttt{\tiny{[GM$|$GM]}}
                \texttt{Relieved: 0} \\
\texttt{Angry: 0} \\
\texttt{Sympathetic: 0} \\
\texttt{Frustrated: 10} \\
            }
        }
    }
    & & \\ \\

    \theutterance \stepcounter{utterance}  
    & & & \multicolumn{2}{p{0.3\linewidth}}{
        \cellcolor[rgb]{0.9,0.9,0.9}{
            \makecell[{{p{\linewidth}}}]{
                \texttt{\tiny{[GM$|$GM]}}
                \texttt{Relieved: 3} \\
\texttt{Angry: 6} \\
\texttt{Sympathetic: 0} \\
\texttt{Frustrated: 8} \\
            }
        }
    }
    & & \\ \\

    \theutterance \stepcounter{utterance}  
    & & & \multicolumn{2}{p{0.3\linewidth}}{
        \cellcolor[rgb]{0.9,0.9,0.9}{
            \makecell[{{p{\linewidth}}}]{
                \texttt{\tiny{[GM$|$GM]}}
                \texttt{game\_result = LOSE} \\
            }
        }
    }
    & & \\ \\

    \theutterance \stepcounter{utterance}  
    & & \multicolumn{4}{p{0.6\linewidth}}{
        \cellcolor[rgb]{0.95,0.95,0.95}{
            \makecell[{{p{\linewidth}}}]{
                \texttt{\tiny{[P1$\leftrightarrow$P]}}
                \texttt{Post{-}episode Reflexion (not sent to GM):} \\
\texttt{New plan: My previous attempts failed because I oscillated between assigning a non{-}zero score to "Angry" and then correcting to zero, indicating a lack of confidence in the emotional nuance. The character's final line, "I'm done here," combined with the description of the opponent's worldview as a "nightmare," signals a definitive end to a futile argument rather than an outburst of rage. The dominant emotion is clearly "Frustrated" (exhaustion with the futility of the debate), while "Angry" is a secondary, less intense reaction to the personal attacks, and "Relieved" and "Sympathetic" are absent. To succeed, I will assign the maximum score (10) to "Frustrated" to reflect the intensity of the exasperation, assign a moderate non{-}zero score (e.g., 4) to "Angry" to acknowledge the sharpness of the final retorts without equating it to the primary emotion, and strictly assign 0 to "Relieved" and "Sympathetic." This balanced approach captures the complexity of the character's state: primarily exhausted and frustrated, but with an undercurrent of irritation.} \\
            }
        }
    }
    & \\ \\

\end{supertabular}
}

\end{document}
